\subsection{Consensus clustering of the learned representation (RQ3)}\label{sec:res_consensus}

With all five autoencoder seeds in the ensemble (\nMembersAe{} members), the consensus Silhouette of the learned representation peaks at $k=\aeAllK$ (\aeAllSil, PAC \aeAllPac) and declines thereafter (\aeAllSilTwelve{} at $k=12$); at $k=8$ only \aeAllStableEight\% of players have a stability score of 0.5 or more (Figure~\ref{fig:consensus}, Table~\ref{tab:baselines}). When the ensemble is restricted to a single autoencoder seed, so that only subsampling varies, the consensus is far tighter (\aeSeedZeroSil{} at $k=3$, \aeSeedZeroSilEight{} at $k=8$, \aeSeedZeroStableEight\% stable). The two ensembles differ only in the inclusion of independently trained encoders, so the dominant source of irreproducibility is the training seed: two runs of the same model on the same data learn latent geometries whose partitions agree only moderately (ARI \ariASeedZero{} between the single-seed and the five-seed consensus). Leaving the latent coordinates unstandardised (\aeNostdSil{} at $k=3$), inserting UMAP before $k$-means (\aeUmapSil{} at $k=3$ but \aeUmapSilEight{} at $k=8$) or embedding HDBSCAN itself in the consensus (PAC \aeHdbPac, \aeHdbStable\% stable) does not change this, and neither does latent dimensionality (coarse consensus Silhouette \aeEightSil{} to \aeOneTwoEightSil{} for $d=8$ to $128$; Table~\ref{tab:sensitivity}a). For the learned representation, a three-group coarse structure, essentially level and length of career, is reproducible across seeds; anything finer is not.

\begin{figure}[t]\centering
\includegraphics[width=\textwidth]{fig5_consensus}
\caption{Consensus diagnostics by number of clusters $k$ for the main representations and two structureless references (left: PAC; centre: consensus Silhouette), and the co-association matrix of the summary-statistic representation ordered by the $k=8$ consensus partition (right; black lines separate the types).}
\label{fig:consensus}
\end{figure}

\subsection{Learned representation versus baselines (RQ4)}\label{sec:res_baselines}

Table~\ref{tab:baselines} compares every representation under the same protocol. The nine summary statistics and the resampled trajectory are markedly more reproducible than any variant of the learned representation: at $k=8$ their consensus Silhouettes are \summSilEight{} and \rawSilEight{} against \aeAllSilEight, with \summStableEight\% and \rawStableEight\% of players stable against \aeAllStableEight\%; they explain more trajectory-shape variance ($R^2_{\mathrm{shape}}$ \summRsqEight{} and \rawRsqEight{} against \aeAllRsqEight) and are at least as strongly associated with All-Star selection and career Win Shares. Ward linkage on the summary statistics (\summWardSilEight) and PCA of the resampled trajectory (\pcaSilEight) are intermediate. DTW $k$-means on the raw sequences, the most direct classical competitor, \dtwSentence

Reproducibility within a representation does not translate into agreement between representations: the $k=8$ partitions of the summary statistics and the resampled trajectory agree at ARI \xSummaryRaw, and every pairing involving the autoencoder lies between \xMin{} and \xMax{} (Appendix Table~\ref{tab:crossrep}). Each representation partitions the same players reproducibly but differently. RQ4 is answered in the negative for the learned representation, and the disagreement raises the question addressed next: whether \emph{any} of these partitions reflects discrete structure.

\begin{table}[t]\centering\scriptsize
\caption{All representations under the same consensus protocol. Sil.$_C$ = consensus Silhouette; PAC = proportion of ambiguous pairs; Stable \% = share of players with stability $\ge 0.5$; $R^2_{\mathrm{shape}}$ = share of resampled-trajectory variance explained by cluster membership; $V_{\mathrm{AS}}$ = Cram\'er's $V$ with All-Star selection; $\varepsilon^2_{\mathrm{WS}}$ = Kruskal--Wallis effect size for career Win Shares. Left block: $k$ selected by the rule; right block: all representations at $k=8$.}
\label{tab:baselines}
\resizebox{\textwidth}{!}{\begin{tabular}{lcccccccccccc}\toprule
 & & \multicolumn{5}{c}{Selected $k$ (rule)} & \multicolumn{6}{c}{Matched comparison at $k=8$} \\ \cmidrule(lr){3-7}\cmidrule(lr){8-13}
Representation & Members & $k$ & Sil.$_C$ & PAC & Stable \% & $R^2_{\mathrm{shape}}$ & Sil.$_C$ & PAC & Stable \% & $R^2_{\mathrm{shape}}$ & $V_{\mathrm{AS}}$ & $\varepsilon^2_{\mathrm{WS}}$ \\ \midrule
Autoencoder $d=64$, 5 seeds (latent space) & 40 & 3 & 0.48 & 0.57 & 80 & 0.37 & 0.44 & 0.32 & 67 & 0.50 & 0.61 & 0.65 \\
Autoencoder $d=64$, seed 0 only & 20 & 3 & 0.90 & 0.17 & 100 & 0.41 & 0.78 & 0.13 & 95 & 0.50 & 0.56 & 0.60 \\
Autoencoder $d=64$, 5 seeds, unstandardised & 40 & 3 & 0.56 & 0.50 & 94 & 0.46 & 0.45 & 0.31 & 70 & 0.57 & 0.63 & 0.67 \\
Autoencoder $d=64$, 5 seeds + UMAP (5-D) & 40 & 3 & 0.71 & 0.38 & 95 & 0.24 & 0.38 & 0.33 & 49 & 0.29 & 0.40 & 0.70 \\
Autoencoder + UMAP + HDBSCAN (consensus) & 40 & 5 & 0.43 & 0.76 & 66 & 0.18 & 0.36 & 0.76 & 66 & 0.18 & 0.29 & 0.54 \\
Autoencoder $d=8$, 3 seeds & 24 & 3 & 0.50 & 0.55 & 90 & 0.40 & 0.34 & 0.35 & 46 & 0.48 & 0.60 & 0.58 \\
Autoencoder $d=16$, 3 seeds & 24 & 4 & 0.57 & 0.41 & 86 & 0.38 & 0.35 & 0.32 & 53 & 0.48 & 0.63 & 0.61 \\
Autoencoder $d=32$, 3 seeds & 24 & 3 & 0.62 & 0.46 & 91 & 0.39 & 0.37 & 0.33 & 50 & 0.49 & 0.62 & 0.64 \\
Autoencoder $d=128$, 3 seeds & 24 & 3 & 0.70 & 0.39 & 92 & 0.39 & 0.49 & 0.29 & 71 & 0.46 & 0.62 & 0.63 \\
Summary statistics, $k$-means & 8 & 4 & 0.93 & 0.08 & 99 & 0.47 & 0.90 & 0.07 & 99 & 0.60 & 0.68 & 0.75 \\
Summary statistics, Ward & 8 & 3 & 0.49 & 0.62 & 96 & 0.42 & 0.47 & 0.33 & 81 & 0.55 & 0.66 & 0.73 \\
Resampled trajectory, $k$-means & 8 & 3 & 0.93 & 0.10 & 99 & 0.57 & 0.85 & 0.10 & 97 & 0.72 & 0.65 & 0.71 \\
PCA of resampled trajectory, $k$-means & 8 & 3 & 0.74 & 0.31 & 94 & 0.14 & 0.50 & 0.24 & 80 & 0.34 & 0.48 & 0.55 \\
DTW $k$-means on sequences & 8 & 3 & 0.76 & 0.34 & 97 & 0.53 & 0.36 & 0.35 & 70 & 0.66 & 0.64 & 0.51 \\
\bottomrule\end{tabular}}
\end{table}

\subsection{Discrete types or a continuum?}\label{sec:res_structure}

Three tests were applied (Table~\ref{tab:structure}). The gap statistic \citep{tibshirani2001} grows monotonically in $k$ against a uniform reference for every representation, the signature of a continuous correlated cloud, and is close to zero for every $k$ against a multivariate-normal reference with the same covariance (maximum \gapSummNormalMax{} for the summary statistics, \gapRawNormalMax{} for the resampled trajectory, \gapAeNormalMax{} for the learned representation); the gap rule selects $k=1$ for all three. Hartigan's dip test \citep{hartigan1985} finds no multimodality on the first three principal components of any representation ($p$ between \dipPmin{} and \dipPmax); among the summary statistics only the number of seasons, a discrete count with a long right tail, rejects unimodality. Finally, the consensus protocol was applied to structureless references \citep{senbabaoglu2014}. \nullSentence

The three tests agree: the reproducible partitions of Section~\ref{sec:res_baselines} are reproducible segmentations of a continuum, not recoveries of discrete career types.

\begin{table}[t]\centering\small
\caption{Tests for discrete structure. Gap statistic: maximum over $k=2,\dots,12$ against a multivariate-normal reference with the sample covariance and the $k$ selected by the gap rule. Dip test: $p$-values of Hartigan's test on the first three principal components. Null consensus: mean consensus Silhouette over three simulated datasets at $k=4$ and $k=8$ for a multivariate-normal reference and for a Gaussian copula with the real marginals. $^a$Gap and dip statistics for the autoencoder are computed on the seed-0 latent space; consensus values are for the five-seed ensemble.}
\label{tab:structure}
\resizebox{\textwidth}{!}{\begin{tabular}{lcccccc}\toprule
 & \multicolumn{2}{c}{Gap statistic (normal ref.)} & Dip test & \multicolumn{3}{c}{Consensus Silhouette, $k=4$ / $k=8$} \\ \cmidrule(lr){2-3}\cmidrule(lr){4-4}\cmidrule(lr){5-7}
Representation & max gap & selected $k$ & min $p$ (PC1--3) & Real data & Gaussian null & Copula null \\ \midrule
Summary statistics & 0.16 & 1 & 0.71 & 0.93 / 0.90 & 0.93 / 0.56 & 0.90 / 0.75 \\
Resampled trajectory & 0.01 & 1 & 0.99 & 0.87 / 0.85 & 0.93 / 0.63 & 0.87 / 0.66 \\
Autoencoder $d=64$, 5 seeds$^a$ & 0.11 & 1 & 0.93 & 0.39 / 0.44 & -- / -- & -- / -- \\
\bottomrule\end{tabular}}
\end{table}

\subsection{A continuous description of careers}\label{sec:res_continuous}

If careers form a continuum, the natural summary is a small set of continuous coordinates. Principal components of the resampled trajectory (ten points of normalised career time, PER units) provide one (Figure~\ref{fig:continuous}). The first component (\pcOneVar\% of variance) is career level: $r=\pcOneMean$ with career-mean PER, \pcOnePeak{} with peak PER, \pcOneWs{} with career Win Shares. The second (\pcTwoVar\%) is the tilt of the career, rising versus declining: uncorrelated with level ($|r|=\pcTwoMean$), $r=\pcTwoSlope$ with the linear slope and \pcTwoPeakpos{} with the relative position of the peak. Career length is a third, separate coordinate. A linear classifier on level, tilt and log length recovers the eight fine types of Section~\ref{sec:profiles} with \ldaAccThree\% cross-validated accuracy (\ldaAccTwo\% from level and tilt alone), so most of the typology is a discretisation of these three coordinates: the types are contiguous bands in the level--tilt plane with no empty regions between them, and the density of players in that plane is a single unimodal cloud.

The same holds when level is removed. On the shape-only representation (each trajectory minus its own mean), the consensus protocol is again highly reproducible ($k=\shapeK$, consensus Silhouette \shapeSil) and returns the early, middle and late peakers of \citet{wakim2014} (Table~\ref{tab:shape}; median peak position \shapePpA, \shapePpB{} and \shapePpC{} of the career). But the dip test finds no multimodality in the shape components ($p\ge\shapeDipPmin$) or in the slope ($p=\dipSlopeP$), and the shape groups differ strongly in career length (median \shapeSeasA, \shapeSeasB{} and \shapeSeasC{} seasons; $\varepsilon^2=\shapeLenEps$), as expected when short careers are stretched onto normalised time. Peak timing is a continuous property of this population: \peakFirstThird\% of players peak in the first third of their career and \peakLastThird\% in the last third, with no gap in between.

\begin{figure}[t]\centering
\includegraphics[width=\textwidth]{fig9_continuous}
\caption{A continuous description. Left: mean resampled PER trajectory and the effect of one standard deviation of the first two principal components (PC1, career level; PC2, rise versus decline). Centre: players in the PC1--PC2 plane coloured by fine type, with selected players labelled. Right: density of players in the same plane.}
\label{fig:continuous}
\end{figure}

\begin{table}[t]\centering\small
\caption{Shape-only consensus groups ($k=\shapeK$; resampled trajectory minus the player's own mean, length excluded), ordered by median peak position; medians per group.}
\label{tab:shape}
\begin{tabular}{lcccccccc}\toprule
Shape group & $n$ & Peak position & Slope & First-season PER & Last-season PER & Mean PER & Seasons & All-Star \% \\ \midrule
S1 & 613 & 0.33 & -0.69 & 14.0 & 8.9 & 12.5 & 8 & 17 \\
S2 & 400 & 0.50 & -0.19 & 11.8 & 9.9 & 13.9 & 11 & 30 \\
S3 & 663 & 0.70 & +0.05 & 10.9 & 11.3 & 11.5 & 4 & 3 \\
\bottomrule\end{tabular}
\end{table}

\subsection{A descriptive typology of the continuum}\label{sec:profiles}

Continuous coordinates are the more faithful summary, but a reproducible segmentation remains useful for description, provided it is not mistaken for a discovery of natural kinds. We therefore also report the consensus partition of the most reproducible representation, the summary statistics, at $k=8$ (Table~\ref{tab:profiles8}, Figure~\ref{fig:profiles}), where the consensus Silhouette is \summSilEight{} and \summStableEight\% of players are stable; the rule-selected $k=\mainK$ solution, in which the eight types nest imperfectly (\nestingPurity\% purity), is given in Appendix~\ref{app:nesting}. Labels follow the rule of Section~\ref{sec:methods}. The long-career groups split by level into a high-PER group (T1: \eightNA{} players, median PER \eightPerA, \eightASA\% All-Stars), a moderate-PER group (T3) and a low-PER group (T6); two medium-length groups differ in direction (T2, early peak then decline; T4, late rise); and the short-career groups are flat-declining (T5), rising (T7) and steeply declining (T8). Median stability is at least \eightStabMin{} in every type. Named players illustrate rather than validate: Michael Jordan, Kobe Bryant, Larry Bird, Stephen Curry and LeBron James are all in \typeJordan, as are Isiah Thomas and Steve Nash; Ralph Sampson is in \typeSampson, Darko Mili\v{c}i\'c in \typeMilicic{} and Adam Morrison in \typeMorrison. The players nearest each type's medoid (Table~\ref{tab:profiles8}) are, by construction, not stars, and represent better what a type typically contains.

\begin{table}[t]\centering\scriptsize
\caption{Fine types ($k=8$) of the summary-statistic consensus, ordered by median career-mean PER and length. Labels follow the fixed rule (level: $>18$ high, 12--18 moderate, $<12$ low; duration: $\ge10$ long, 5--9 medium, $\le4$ short; slope qualifier if $|\text{slope}|>0.25$ PER per season).}
\label{tab:profiles8}
\resizebox{\textwidth}{!}{\begin{tabular}{lcccccccclp{3.6cm}}\toprule
Type & $n$ & Stability & Mean PER & Peak PER & Seasons & Slope & All-Star \% & Career WS & Label (rule) & Nearest to medoid \\ \midrule
T1 & 138 & 0.98 & 18.5 & 23.6 & 14 & -0.23 & 86 & 89.6 & high-PER, long career & Shawn Kemp, Paul George, Sam Cassell \\
T2 & 166 & 0.95 & 15.1 & 18.1 & 6 & -1.02 & 19 & 16.8 & moderate-PER, medium career, declining & Lawrence Funderburke, Edgar Jones, Lewis Lloyd \\
T3 & 283 & 0.96 & 14.2 & 18.4 & 13 & -0.30 & 28 & 49.1 & moderate-PER, long career, declining & Brendan Haywood, Brad Davis, Jeff Malone \\
T4 & 141 & 0.96 & 13.8 & 16.9 & 7 & +0.25 & 6 & 13.4 & moderate-PER, medium career, rising & Marty Conlon, Will Bynum, Darius Miles \\
T5 & 276 & 0.95 & 11.2 & 13.0 & 4 & -0.40 & 0 & 3.8 & low-PER, short career, declining & Bob Wilkerson, Obinna Ekezie, Tom Garrick \\
T6 & 339 & 0.95 & 11.2 & 14.7 & 10 & -0.37 & 1 & 17.1 & low-PER, long career, declining & Dante Cunningham, Brad Lohaus, Kevin Duckworth \\
T7 & 206 & 0.98 & 9.5 & 11.8 & 3 & +0.90 & 0 & 0.8 & low-PER, short career, rising & Scott Roth, Delaney Rudd, Laron Profit \\
T8 & 127 & 0.90 & 9.5 & 12.4 & 3 & -2.90 & 0 & 1.3 & low-PER, short career, declining & Jon Brockman, Chris Welp, Samardo Samuels \\
\bottomrule\end{tabular}}
\end{table}

\begin{figure}[t]\centering
\includegraphics[width=\textwidth]{fig6_profiles}
\caption{PER trajectories of the eight fine types on the career-season axis: individual careers (thin lines, up to 150 per panel), median (black) and interquartile band (shaded), computed at seasons with at least ten observations; $n$ and median stability per type.}
\label{fig:profiles}
\end{figure}

\subsection{External validation (RQ5)}\label{sec:res_external}

Table~\ref{tab:external} relates the eight types to variables that played no part in clustering. Associations with honours and value are large: Cram\'er's $V$ is \VAllStar{} for ever being an All-Star and \VAllNba{} for All-NBA selection; $\varepsilon^2=\epsWs$ for career Win Shares and \epsBpm{} for mean BPM; \eightASA\% of T1 players were All-Stars against at most \eightASmaxOther\% in any other type. The association with draft category is moderate ($V=\VDraft$): \eightLotteryA\% of T1 were lottery picks, whereas the short-career types are dominated by second-round and undrafted players. Position is unrelated to type ($V=\VPos$, $p=\pPos$): the types describe productivity paths, not roles, consistent with \citet{wakim2014}. Because All-Star selection, WS and BPM share inputs with PER, these associations show that the types are not artefacts of the representation, not that they carry information beyond box-score productivity.

\begin{table}[t]\centering\scriptsize
\caption{External variables by fine type. Draft category and position are row percentages; honours are the share of players with at least one selection; WS, BPM and games are medians. Last row: effect sizes over the eight types.}
\label{tab:external}
\resizebox{\textwidth}{!}{\begin{tabular}{lccccccccccc}\toprule
 & \multicolumn{4}{c}{Draft category (\% of type)} & Position & \multicolumn{2}{c}{Honours (\%)} & \multicolumn{3}{c}{Career value (median)} \\ \cmidrule(lr){2-5}\cmidrule(lr){6-6}\cmidrule(lr){7-8}\cmidrule(lr){9-11}
Type & Lottery & 1st rd & 2nd rd & Undrafted & PG/SG/SF/PF/C & All-Star & All-NBA & WS & BPM & Games \\ \midrule
T1 & 71 & 15 & 9 & 5 & 17/15/17/25/26 & 86 & 68 & 89.6 & 1.8 & 968 \\
T2 & 28 & 16 & 17 & 39 & 14/15/22/27/22 & 19 & 6 & 16.8 & -0.7 & 357 \\
T3 & 47 & 29 & 16 & 7 & 25/18/19/22/17 & 28 & 9 & 49.1 & -0.3 & 866 \\
T4 & 26 & 37 & 21 & 16 & 21/18/19/16/25 & 6 & 2 & 13.4 & -1.1 & 396 \\
T5 & 13 & 25 & 33 & 29 & 20/17/19/23/21 & 0 & 0 & 3.8 & -2.7 & 216 \\
T6 & 22 & 31 & 29 & 18 & 19/20/21/19/20 & 1 & 0 & 17.1 & -2.2 & 584 \\
T7 & 12 & 25 & 43 & 20 & 22/27/17/15/20 & 0 & 0 & 0.8 & -3.7 & 134 \\
T8 & 12 & 22 & 32 & 34 & 18/20/24/17/21 & 0 & 0 & 1.3 & -3.9 & 106 \\
\midrule Cram\'er's $V$ / $\varepsilon^2$ & \multicolumn{4}{c}{$V=0.27$} & $V=0.08$ & $V=0.68$ & $V=0.68$ & $\varepsilon^2=0.75$ & $\varepsilon^2=0.54$ & $\varepsilon^2=0.74$ \\
\bottomrule\end{tabular}}
\end{table}

\subsection{Sensitivity analyses (RQ5)}\label{sec:sensitivity}

Table~\ref{tab:sensitivity} reports agreement with the reference partition when one element of the analysis is changed. For the summary-statistic typology (panel b), dropping the games rule and deleting gaps leave the $k=8$ partition essentially unchanged (ARI \ariSEightSummaryNogames{} and \ariSEightSummaryConsecutive), whereas restrictions that remove players change it more (ARI \ariSEightMinRestrict{} to \ariSEightMaxRestrict), partly because the rule then selects $k=3$ rather than $4$ and partly because $k$-means boundaries move when the sample distribution is reshaped. Replacing PER by BPM changes the partition substantially (ARI \ariSEightBpm) although career-mean PER and BPM correlate at $r=\corrPerBpm$: the segmentation is metric-specific. Both patterns are what a continuum implies. For the learned representation (panel a) every change, including changes that alter nothing but the training seed, moves the partition substantially (ARI \ariAMin{} to \ariAMax). Finally, the nineteen-cluster solution of the earlier version, matched by name for \nMatchedOld{} clustered players, has an ARI of \ariOldEight{} with the fine typology and \ariOldAe{} with the learned-representation consensus; it is not recovered by any reproducible analysis.

\begin{table}[t]\centering\small
\caption{Sensitivity of the partitions: ARI with the reference partition on players common to both analyses, at the selected $k$ of each variant and at $k=8$.}
\label{tab:sensitivity}
\resizebox{\textwidth}{!}{\begin{tabular}{lcccc}\toprule
Variant & Common players & Selected $k$ & ARI with reference (selected $k$) & ARI with reference ($k=8$) \\ \midrule
\multicolumn{5}{l}{\textit{(a) Autoencoder pipeline, reference: AE $d$=64, five seeds, $k$-means consensus}} \\
Single AE seed & 1676 & 3 & 0.58 & 0.57 \\
Unstandardised latent & 1676 & 3 & 0.54 & 0.52 \\
UMAP before $k$-means & 1676 & 3 & 0.31 & 0.29 \\
UMAP + HDBSCAN & 1676 & 5 & 0.24 & 0.18 \\
AE $d$=8 & 1676 & 3 & 0.33 & 0.38 \\
AE $d$=16 & 1676 & 4 & 0.29 & 0.40 \\
AE $d$=32 & 1676 & 3 & 0.54 & 0.51 \\
AE $d$=128 & 1676 & 3 & 0.62 & 0.53 \\
2025 restrictions & 1676 & 3 & 0.20 & 0.18 \\
250-minute threshold & 1544 & 3 & 0.35 & 0.29 \\
$\ge$3 seasons & 1444 & 3 & 0.35 & 0.27 \\
No games rule & 1676 & 3 & 0.41 & 0.38 \\
Debut $\le$2005--06 & 1437 & 3 & 0.39 & 0.37 \\
Gaps deleted & 1676 & 3 & 0.62 & 0.38 \\
BPM instead of PER & 1676 & 3 & 0.32 & 0.24 \\
\midrule \multicolumn{5}{l}{\textit{(b) Summary-statistic typology, reference: primary consensus}} \\
2025 restrictions & 1676 & 3 & 0.31 & 0.51 \\
250-minute threshold & 1544 & 3 & 0.48 & 0.59 \\
$\ge$3 seasons & 1444 & 3 & 0.44 & 0.74 \\
No games rule & 1676 & 4 & 0.92 & 0.92 \\
Debut $\le$2005--06 & 1437 & 3 & 0.41 & 0.82 \\
Gaps deleted & 1676 & 4 & 0.92 & 0.88 \\
BPM instead of PER & 1676 & 3 & 0.29 & 0.27 \\
Ward instead of $k$-means & 1676 & 3 & 0.33 & 0.53 \\
\bottomrule\end{tabular}}
\end{table}
